\documentclass[a4paper,12pt]{ltjsarticle}

% LuaLaTeX用日本語設定
\usepackage{luatexja}
\usepackage[ipaex]{luatexja-preset}

% 数学パッケージ
\usepackage{amsmath, amssymb, amsthm}
\usepackage{mathtools}

% 可換図式
\usepackage{tikz-cd}
\usepackage{tikz}
\usetikzlibrary{positioning}

% その他のパッケージ
\usepackage{enumerate}
\usepackage{url}
\usepackage{hyperref}

% 定理環境
\theoremstyle{definition}
\newtheorem{definition}{定義}[section]
\newtheorem{theorem}[definition]{定理}
\newtheorem{proposition}[definition]{命題}
\newtheorem{lemma}[definition]{補題}
\newtheorem{corollary}[definition]{系}
\newtheorem{example}[definition]{例}
\newtheorem{remark}[definition]{注意}

% タイトル情報
\title{構造塔の圏論的構造\\
\large Leanによる形式化と数学的解説}
\author{%
  TeX文書生成: Claude Code (Anthropic Claude Sonnet 4.5)\\
  Lean形式化: Codex (OpenAI)
}
\date{\today}

\begin{document}

\maketitle

\begin{abstract}
本稿では、Bourbakiの構造理論に基づく\textbf{構造塔}（Structure Tower）の概念を圏論的に定式化し、Lean証明支援系を用いて形式化する。構造塔は集合に階層的な構造を与える枠組みであり、その射と合成が圏をなすことを示す。さらに、普遍性の概念を通じて、自由構造塔や直積の圏論的性質を解明する。最小層を導入することで、完全な普遍性が証明可能となることを示す。
\end{abstract}

\tableofcontents
\newpage

%=============================================================================
\section{序論}
%=============================================================================

構造塔（Structure Tower）は、Nicolas Bourbakiの構造理論において重要な役割を果たす概念である。集合 $X$ に対して、添字集合 $I$ で parameterized された部分集合の族 $(X_i)_{i \in I}$ を考え、これらが特定の性質（被覆性、単調性）を満たすとき、三つ組 $(X, (X_i)_{i \in I}, \leq)$ を構造塔と呼ぶ。

本稿の目的は以下の通りである：
\begin{enumerate}
\item 構造塔の射を定義し、構造塔が圏をなすことを証明する
\item 忘却関手、自由構造塔、同型射などの圏論的構成を与える
\item 普遍性の定式化における数学的問題を明らかにし、解決策を提示する
\item 最小層を持つ構造塔において、完全な普遍性を証明する
\end{enumerate}

本稿の内容は、Lean証明支援系（Mathlib）を用いて形式的に検証されている。

%=============================================================================
\section{基本的な構造塔}
\label{sec:cat1}
%=============================================================================

本節では、構造塔の基本的な定義と、それが圏をなすことを示す。

%-----------------------------------------------------------------------------
\subsection{構造塔の定義}
%-----------------------------------------------------------------------------

\begin{definition}[構造塔]
\label{def:structure-tower}
\textbf{構造塔}（Structure Tower）とは、以下のデータからなる：
\begin{itemize}
\item 基礎集合 $X$（carrier）
\item 添字集合 $I$（Index）
\item $I$ 上の前順序 $\leq$
\item 各層を定める写像 $\mathrm{layer}: I \to \mathcal{P}(X)$
\end{itemize}
これらは以下の公理を満たす：
\begin{enumerate}
\item \textbf{被覆性}（covering）: $\forall x \in X,\ \exists i \in I,\ x \in X_i$
\item \textbf{単調性}（monotone）: $\forall i, j \in I,\ i \leq j \Rightarrow X_i \subseteq X_j$
\end{enumerate}
ここで $X_i := \mathrm{layer}(i)$ とする。
\end{definition}

被覆性は、すべての要素が少なくとも1つの層に属することを保証する。単調性は、順序関係が層の包含関係を保存することを意味する。

%-----------------------------------------------------------------------------
\subsection{構造塔の射}
%-----------------------------------------------------------------------------

\begin{definition}[構造塔の射]
\label{def:tower-morphism}
構造塔 $T = (X, I, \leq, \mathrm{layer})$ と $T' = (X', I', \leq', \mathrm{layer}')$ の間の\textbf{射}（morphism）とは、対 $(f, \varphi)$ であって、以下を満たすもの：
\begin{itemize}
\item $f: X \to X'$ は基礎集合の写像
\item $\varphi: I \to I'$ は順序を保つ写像
\item $\forall i, j \in I,\ i \leq j \Rightarrow \varphi(i) \leq' \varphi(j)$（順序保存）
\item $\forall i \in I,\ \forall x \in X,\ x \in X_i \Rightarrow f(x) \in X'_{\varphi(i)}$（層構造保存）
\end{itemize}
\end{definition}

層構造保存の条件は、射が構造塔の階層構造を尊重することを要求する。

\begin{figure}[h]
\centering
\begin{tikzcd}
X \arrow[r, "f"] \arrow[d, "\text{層への包含}"] & X' \arrow[d, "\text{層への包含}"] \\
X_i \arrow[r, "f|_{X_i}"] & X'_{\varphi(i)}
\end{tikzcd}
\caption{構造塔の射の可換性}
\label{fig:morphism-commute}
\end{figure}

%-----------------------------------------------------------------------------
\subsection{射の合成と恒等射}
%-----------------------------------------------------------------------------

\begin{definition}[射の合成]
\label{def:morphism-composition}
射 $(f, \varphi): T \to T'$ と $(g, \psi): T' \to T''$ の合成は、
\[
(g, \psi) \circ (f, \varphi) := (g \circ f,\ \psi \circ \varphi)
\]
で定義される。
\end{definition}

\begin{definition}[恒等射]
構造塔 $T$ 上の恒等射は、
\[
\mathrm{id}_T := (\mathrm{id}_X, \mathrm{id}_I)
\]
で定義される。
\end{definition}

\begin{lemma}[射の合成の性質]
射の合成は以下の性質を満たす：
\begin{enumerate}
\item \textbf{結合律}: $(h \circ g) \circ f = h \circ (g \circ f)$
\item \textbf{左単位律}: $\mathrm{id}_{T'} \circ f = f$
\item \textbf{右単位律}: $f \circ \mathrm{id}_T = f$
\end{enumerate}
\end{lemma}

%-----------------------------------------------------------------------------
\subsection{構造塔の圏}
%-----------------------------------------------------------------------------

\begin{theorem}[構造塔は圏をなす]
\label{thm:tower-category}
構造塔を対象とし、上記の射を射とすることで、圏 $\mathbf{StructureTower}$ が定義される。
\end{theorem}

\begin{proof}
圏の公理（恒等射の存在、合成の結合律、単位律）は、上記の補題で示した性質により満たされる。
\end{proof}

%-----------------------------------------------------------------------------
\subsection{具体例：自然数の初期区間}
%-----------------------------------------------------------------------------

\begin{example}[自然数上の構造塔]
\label{ex:nat-tower}
以下のように定義される構造塔を考える：
\begin{itemize}
\item 基礎集合: $X = \mathbb{N}$
\item 添字集合: $I = \mathbb{N}$（通常の順序）
\item 層: $X_n = \{k \in \mathbb{N} \mid k \leq n\}$（初期区間）
\end{itemize}

被覆性：任意の $x \in \mathbb{N}$ に対して $x \in X_x$。\\
単調性：$n \leq m$ ならば $X_n \subseteq X_m$。
\end{example}

%-----------------------------------------------------------------------------
\subsection{直積構造塔}
%-----------------------------------------------------------------------------

\begin{definition}[直積構造塔]
\label{def:product-tower}
二つの構造塔 $T_1 = (X_1, I_1, \leq_1, \mathrm{layer}_1)$ と $T_2 = (X_2, I_2, \leq_2, \mathrm{layer}_2)$ の直積は、
\[
T_1 \times T_2 := (X_1 \times X_2,\ I_1 \times I_2,\ \leq_{I_1 \times I_2},\ \mathrm{layer}_{(i,j)})
\]
ここで、
\begin{itemize}
\item $(i_1, i_2) \leq (j_1, j_2) \Leftrightarrow i_1 \leq_1 j_1 \land i_2 \leq_2 j_2$（直積順序）
\item $\mathrm{layer}_{(i,j)} := \mathrm{layer}_1(i) \times \mathrm{layer}_2(j)$
\end{itemize}
\end{definition}

\begin{proposition}
直積構造塔は構造塔の公理（被覆性、単調性）を満たす。
\end{proposition}

%=============================================================================
\section{関手と自然変換}
\label{sec:cat2-advanced}
%=============================================================================

本節では、構造塔の圏から他の圏への関手を定義し、圏論的な構成を探求する。

%-----------------------------------------------------------------------------
\subsection{忘却関手}
%-----------------------------------------------------------------------------

\begin{definition}[基礎集合への忘却関手]
\label{def:forget-carrier}
関手 $U: \mathbf{StructureTower} \to \mathbf{Type}$ を以下のように定義する：
\begin{itemize}
\item 対象への作用: $U(T) := X$（基礎集合）
\item 射への作用: $U(f, \varphi) := f$（基礎写像）
\end{itemize}
\end{definition}

\begin{proposition}
$U$ は関手の公理を満たす：
\begin{enumerate}
\item $U(\mathrm{id}_T) = \mathrm{id}_{U(T)}$
\item $U(g \circ f) = U(g) \circ U(f)$
\end{enumerate}
\end{proposition}

同様に、添字集合への忘却関手 $V: \mathbf{StructureTower} \to \mathbf{Type}$ も定義できる。

%-----------------------------------------------------------------------------
\subsection{同型射}
%-----------------------------------------------------------------------------

\begin{definition}[構造塔の同型]
\label{def:tower-isomorphism}
射 $(f, \varphi): T \to T'$ が\textbf{同型射}であるとは、ある射 $(g, \psi): T' \to T$ が存在して、
\[
(g, \psi) \circ (f, \varphi) = \mathrm{id}_T, \quad (f, \varphi) \circ (g, \psi) = \mathrm{id}_{T'}
\]
を満たすことをいう。
\end{definition}

\begin{proposition}
同型射 $(f, \varphi): T \to T'$ において、基礎写像 $f$ は全単射である。
\end{proposition}

\begin{proof}
逆射 $(g, \psi)$ が存在するので、$g \circ f = \mathrm{id}_X$ かつ $f \circ g = \mathrm{id}_{X'}$ より、$f$ は全単射。
\end{proof}

同型射は以下の性質を持つ：
\begin{itemize}
\item \textbf{反射律}: $\mathrm{id}_T$ は同型射
\item \textbf{対称律}: $(f, \varphi)$ が同型なら $(g, \psi)$ も同型
\item \textbf{推移律}: 同型射の合成は同型射
\end{itemize}

%-----------------------------------------------------------------------------
\subsection{自由構造塔}
%-----------------------------------------------------------------------------

\begin{definition}[自由構造塔（離散版）]
\label{def:free-tower-discrete}
集合 $X$ に対して、\textbf{自由構造塔} $F(X)$ を以下のように定義する：
\begin{itemize}
\item 基礎集合: $X$
\item 添字集合: $X$（離散順序: $i \leq j \Leftrightarrow i = j$）
\item 層: $X_i = \{i\}$（単元集合）
\end{itemize}
\end{definition}

自由構造塔は、各要素が独自の層を持つ最も単純な構造塔である。

%-----------------------------------------------------------------------------
\subsection{直積の射影と普遍性}
%-----------------------------------------------------------------------------

直積構造塔 $T_1 \times T_2$ に対して、射影射を定義できる。

\begin{definition}[射影射]
\label{def:projections}
\begin{align*}
\pi_1 &: T_1 \times T_2 \to T_1, \quad \pi_1(x_1, x_2) = x_1, \quad \pi_1(i_1, i_2) = i_1 \\
\pi_2 &: T_1 \times T_2 \to T_2, \quad \pi_2(x_1, x_2) = x_2, \quad \pi_2(i_1, i_2) = i_2
\end{align*}
\end{definition}

\begin{figure}[h]
\centering
\begin{tikzcd}
& T \arrow[ld, "f_1"'] \arrow[rd, "f_2"] \arrow[d, dashed, "\exists ! h"] & \\
T_1 & T_1 \times T_2 \arrow[l, "\pi_1"] \arrow[r, "\pi_2"'] & T_2
\end{tikzcd}
\caption{直積の普遍性}
\label{fig:product-universality}
\end{figure}

\begin{theorem}[直積の普遍性（弱形式）]
\label{thm:product-universality-weak}
任意の構造塔 $T$ と射 $f_1: T \to T_1$, $f_2: T \to T_2$ に対して、射 $h: T \to T_1 \times T_2$ が存在して、
\[
\pi_1 \circ h = f_1, \quad \pi_2 \circ h = f_2
\]
を満たす。ただし、$h$ の一意性は基礎写像に関してのみ成立する。
\end{theorem}

この段階では、添字写像の選択の自由度により、完全な一意性は保証されない。

%=============================================================================
\section{普遍性の問題と解決}
\label{sec:cat2-revised}
%=============================================================================

本節では、構造塔の普遍性に関する数学的問題を明らかにし、複数のアプローチを比較する。

%-----------------------------------------------------------------------------
\subsection{問題の明確化}
%-----------------------------------------------------------------------------

\begin{remark}[普遍性の非一意性]
\label{rem:non-uniqueness-issue}
構造塔の定義（定義\ref{def:structure-tower}）において、被覆性の公理
\[
\forall x \in X,\ \exists i \in I,\ x \in X_i
\]
は、各要素が\textbf{少なくとも1つ}の層に属することを保証するが、\textbf{複数の層}に属することも許される。

したがって、射 $(f, \varphi): T \to T'$ を定義する際、同じ基礎写像 $f$ に対して異なる添字写像 $\varphi$ を持つ複数の射が存在しうる。これにより、普遍性における完全な一意性 $\exists!$ が成立しない。
\end{remark}

\begin{example}[一意性の破れの例]
実数区間の構造塔を考える：
\begin{itemize}
\item 基礎集合: $X = [0, 1]$
\item 添字集合: $I = [0, 1]$
\item 層: $X_t = [0, t]$
\end{itemize}

任意の点 $x \in (0, 1)$ は無限個の層 $X_t$（$t \geq x$）に属する。したがって、ある写像 $f: X \to X'$ に対して、異なる添字写像を持つ複数の射が構成可能である。
\end{example}

%-----------------------------------------------------------------------------
\subsection{三つのアプローチ}
%-----------------------------------------------------------------------------

この問題に対して、以下の三つのアプローチが考えられる：

\subsubsection{Version A: 最小層を持つ構造塔}

各要素が属する\textbf{最小の層}を選ぶ関数を追加する。

\begin{definition}[最小層を持つ構造塔]
\label{def:tower-with-min}
\textbf{最小層を持つ構造塔} $T$ は、通常の構造塔のデータに加えて：
\begin{itemize}
\item 関数 $\mathrm{minLayer}: X \to I$
\end{itemize}
を持ち、以下を満たす：
\begin{enumerate}
\item $\forall x \in X,\ x \in X_{\mathrm{minLayer}(x)}$（包含性）
\item $\forall x \in X,\ \forall i \in I,\ x \in X_i \Rightarrow \mathrm{minLayer}(x) \leq i$（最小性）
\end{enumerate}
\end{definition}

このとき、射の定義に追加条件を課す。

\begin{definition}[最小層を保存する射]
射 $(f, \varphi): T \to T'$ が\textbf{最小層を保存する}とは、
\[
\forall x \in X,\ \mathrm{minLayer}'(f(x)) = \varphi(\mathrm{minLayer}(x))
\]
を満たすことをいう。
\end{definition}

この条件により、添字写像 $\varphi$ が一意に決定される。

\subsubsection{Version B: 基礎写像のみの一意性}

元の定義を保ち、普遍性を「基礎写像の一意性」に弱める。

\begin{definition}[弱普遍性]
\label{def:weak-universality}
構成 $F$ が\textbf{弱普遍性}を持つとは、任意の写像 $f: X \to Y$ に対して、射 $(h, \psi)$ が存在して、
\[
U(h, \psi) = f
\]
を満たし、さらに $h$ が一意であることをいう（ただし $\psi$ は一意とは限らない）。
\end{definition}

\subsubsection{Version C: 分離された構造塔}

各要素が正確に1つの層にのみ属するという強い条件を課す。

\begin{definition}[分離構造塔]
\label{def:separated-tower}
構造塔 $T$ が\textbf{分離されている}とは、
\[
\forall i, j \in I,\ i \neq j \Rightarrow X_i \cap X_j = \emptyset
\]
を満たすことをいう。
\end{definition}

この定義は非常に制限的であり、多くの自然な例（連続的な構造塔）を除外する。

%-----------------------------------------------------------------------------
\subsection{アプローチの比較}
%-----------------------------------------------------------------------------

\begin{center}
\begin{tabular}{|l|c|c|c|}
\hline
\textbf{性質} & \textbf{Version A} & \textbf{Version B} & \textbf{Version C} \\
\hline
完全な一意性 & \checkmark & $\times$ & \checkmark \\
一般性 & 高い & 最も高い & 低い \\
証明の容易さ & 中程度 & やや難 & 容易 \\
実用性 & 高い & 高い & 低い \\
\hline
\end{tabular}
\end{center}

\begin{remark}[推奨事項]
\begin{itemize}
\item \textbf{学習目的}: Version A を推奨。完全な普遍性が証明可能で、概念的にも自然。
\item \textbf{研究・一般理論}: Version B を推奨。より一般的な設定で、「本質的な」部分（基礎写像）の一意性が保証される。
\item \textbf{特殊な応用}: Version C は制限的だが、単純な証明が可能。
\end{itemize}
\end{remark}

%=============================================================================
\section{完全な形式化}
\label{sec:cat2-complete}
%=============================================================================

本節では、Version A（最小層を持つ構造塔）の完全な形式化を与え、普遍性の完全な証明を提示する。

%-----------------------------------------------------------------------------
\subsection{最小層を持つ構造塔の圏}
%-----------------------------------------------------------------------------

最小層を持つ構造塔（定義\ref{def:tower-with-min}）と最小層を保存する射により、圏 $\mathbf{StructureTowerWithMin}$ が定義される。

\begin{theorem}
$\mathbf{StructureTowerWithMin}$ は圏の公理を満たす。
\end{theorem}

\begin{proof}
\begin{enumerate}
\item \textbf{恒等射}: $\mathrm{id}_T = (\mathrm{id}_X, \mathrm{id}_I)$ は明らかに最小層を保存する。
\item \textbf{合成の well-defined 性}: $(g, \psi) \circ (f, \varphi)$ が最小層を保存することを示す：
\begin{align*}
\mathrm{minLayer}''(g(f(x))) &= \psi(\mathrm{minLayer}'(f(x))) \quad (\text{$g$ の性質}) \\
&= \psi(\varphi(\mathrm{minLayer}(x))) \quad (\text{$f$ の性質})
\end{align*}
これは合成 $\psi \circ \varphi$ が最小層を保存することを示す。
\item 結合律、単位律は通常の関数合成の性質から従う。
\end{enumerate}
\end{proof}

%-----------------------------------------------------------------------------
\subsection{自由構造塔の普遍性}
%-----------------------------------------------------------------------------

\begin{definition}[半順序集合上の自由構造塔]
\label{def:free-tower-min}
半順序集合 $(X, \leq)$ に対して、自由構造塔 $F_{\min}(X)$ を以下のように定義する：
\begin{itemize}
\item 基礎集合: $X$
\item 添字集合: $X$（元の半順序）
\item 層: $X_i = \{x \in X \mid x \leq i\}$（下閉集合）
\item 最小層: $\mathrm{minLayer}(x) = x$
\end{itemize}
\end{definition}

\begin{lemma}
$F_{\min}(X)$ において、$\mathrm{minLayer}$ は最小性を満たす。
\end{lemma}

\begin{proof}
$x \in X_i$ すなわち $x \leq i$ ならば、$\mathrm{minLayer}(x) = x \leq i$。
\end{proof}

\begin{theorem}[自由構造塔の普遍性]
\label{thm:free-tower-universality}
半順序集合 $X$ と構造塔 $T$ に対して、写像 $f: X \to \mathrm{carrier}(T)$ が
\[
\text{$x \mapsto \mathrm{minLayer}_T(f(x))$ が単調}
\]
を満たすならば、一意的な射 $\varphi: F_{\min}(X) \to T$ が存在して、
\[
U(\varphi) = f
\]
を満たす（ここで $U$ は基礎集合への忘却関手）。
\end{theorem}

\begin{proof}
\textbf{（存在性）}
射 $\varphi = (f, \psi)$ を以下のように定義する：
\begin{itemize}
\item $\varphi_{\text{map}} := f$
\item $\varphi_{\text{indexMap}}(x) := \mathrm{minLayer}_T(f(x))$
\end{itemize}

\begin{enumerate}
\item \textbf{順序保存}: $x \leq y$ ならば、仮定より $\mathrm{minLayer}_T(f(x)) \leq \mathrm{minLayer}_T(f(y))$。
\item \textbf{層構造保存}: $x \in X_i$ すなわち $x \leq i$ とする。このとき、
\begin{align*}
\mathrm{minLayer}_T(f(x)) &\leq \mathrm{minLayer}_T(f(i)) \quad (\text{単調性})
\end{align*}
したがって、
\[
f(x) \in T_{\mathrm{minLayer}_T(f(x))} \subseteq T_{\mathrm{minLayer}_T(f(i))} = T_{\varphi_{\text{indexMap}}(i)}
\]
\item \textbf{最小層保存}: 定義より、
\[
\mathrm{minLayer}_T(f(x)) = \varphi_{\text{indexMap}}(\mathrm{minLayer}_{F_{\min}(X)}(x)) = \varphi_{\text{indexMap}}(x)
\]
\end{enumerate}

\textbf{（一意性）}
射 $\psi = (g, \chi)$ が $U(\psi) = f$ を満たすとする。このとき、
\begin{itemize}
\item $g = f$（基礎写像）
\item 最小層保存より、
\[
\chi(x) = \chi(\mathrm{minLayer}_{F_{\min}(X)}(x)) = \mathrm{minLayer}_T(g(x)) = \mathrm{minLayer}_T(f(x)) = \varphi_{\text{indexMap}}(x)
\]
\end{itemize}
したがって、$\psi = \varphi$。
\end{proof}

この定理により、最小層を持つ構造塔の枠組みでは、\textbf{完全な普遍性}（$\exists!$）が証明可能となる。

%-----------------------------------------------------------------------------
\subsection{直積の普遍性}
%-----------------------------------------------------------------------------

\begin{definition}[最小層を持つ直積構造塔]
$T_1, T_2$ の直積 $T_1 \times T_2$ において、
\[
\mathrm{minLayer}_{T_1 \times T_2}(x_1, x_2) := (\mathrm{minLayer}_{T_1}(x_1), \mathrm{minLayer}_{T_2}(x_2))
\]
と定義する。
\end{definition}

\begin{theorem}[直積の普遍性]
\label{thm:product-universality-complete}
任意の構造塔 $T$ と射 $f_1: T \to T_1$, $f_2: T \to T_2$ に対して、一意的な射 $h: T \to T_1 \times T_2$ が存在して、
\[
\pi_1 \circ h = f_1, \quad \pi_2 \circ h = f_2
\]
を満たす。
\end{theorem}

\begin{proof}
\textbf{（存在性）}
$h = (h_{\text{map}}, h_{\text{indexMap}})$ を以下のように定義：
\begin{align*}
h_{\text{map}}(x) &:= (f_{1,\text{map}}(x), f_{2,\text{map}}(x)) \\
h_{\text{indexMap}}(i) &:= (f_{1,\text{indexMap}}(i), f_{2,\text{indexMap}}(i))
\end{align*}

最小層保存の確認：
\begin{align*}
\mathrm{minLayer}_{T_1 \times T_2}(h_{\text{map}}(x)) &= (\mathrm{minLayer}_{T_1}(f_{1,\text{map}}(x)), \mathrm{minLayer}_{T_2}(f_{2,\text{map}}(x))) \\
&= (f_{1,\text{indexMap}}(\mathrm{minLayer}_T(x)), f_{2,\text{indexMap}}(\mathrm{minLayer}_T(x))) \\
&= h_{\text{indexMap}}(\mathrm{minLayer}_T(x))
\end{align*}

\textbf{（一意性）}
射 $k$ が同じ条件を満たすとする。射影との可換性より、
\begin{align*}
k_{\text{map}}(x) &= (f_{1,\text{map}}(x), f_{2,\text{map}}(x)) = h_{\text{map}}(x) \\
k_{\text{indexMap}}(i) &= (f_{1,\text{indexMap}}(i), f_{2,\text{indexMap}}(i)) = h_{\text{indexMap}}(i)
\end{align*}
したがって、$k = h$。
\end{proof}

\begin{figure}[h]
\centering
\begin{tikzcd}
& T \arrow[ld, "f_1"'] \arrow[rd, "f_2"] \arrow[d, dashed, "\exists ! h"] & \\
T_1 & T_1 \times T_2 \arrow[l, "\pi_1"] \arrow[r, "\pi_2"'] & T_2 \\
& \text{最小層保存により} \\
& h_{\text{indexMap}} \text{ が一意に決まる}
\end{tikzcd}
\caption{最小層を持つ直積の完全な普遍性}
\label{fig:product-universality-complete}
\end{figure}

%-----------------------------------------------------------------------------
\subsection{具体例：自然数の構造塔}
%-----------------------------------------------------------------------------

\begin{example}[後者関数が誘導する射の一意性]
自然数上の自由構造塔 $F_{\min}(\mathbb{N})$ を考える。後者関数 $\mathrm{succ}: \mathbb{N} \to \mathbb{N}$ は単調なので、定理\ref{thm:free-tower-universality}により、一意的な射
\[
\varphi: F_{\min}(\mathbb{N}) \to F_{\min}(\mathbb{N})
\]
が存在して、$U(\varphi) = \mathrm{succ}$ を満たす。

具体的には、
\begin{align*}
\varphi_{\text{map}}(n) &= n + 1 \\
\varphi_{\text{indexMap}}(n) &= \mathrm{minLayer}(n+1) = n + 1
\end{align*}
となる。
\end{example}

%=============================================================================
\section{結論と展望}
%=============================================================================

本稿では、構造塔の圏論的構造を Lean 証明支援系を用いて形式化した。主な成果は以下の通りである：

\begin{enumerate}
\item \textbf{基本的な圏構造}（第\ref{sec:cat1}節）：構造塔の射、合成、恒等射を定義し、圏をなすことを証明した。

\item \textbf{圏論的構成}（第\ref{sec:cat2-advanced}節）：忘却関手、同型射、自由構造塔、直積などの基本的な構成を与えた。

\item \textbf{普遍性の問題の解明}（第\ref{sec:cat2-revised}節）：被覆性の公理から生じる添字写像の非一意性を明らかにし、三つの解決アプローチを提示した。

\item \textbf{完全な形式化}（第\ref{sec:cat2-complete}節）：最小層を持つ構造塔において、自由構造塔と直積の完全な普遍性（存在性と一意性）を証明した。
\end{enumerate}

\subsection{数学的洞察}

本研究を通じて得られた重要な洞察は以下の通りである：

\begin{itemize}
\item \textbf{形式化における隠れた仮定}：非形式的な数学では「自然に選べる」と思われる操作（層の選択）が、形式的には選択の自由度を持つことがある。

\item \textbf{普遍性の階層}：「射としての一意性」と「基礎写像としての一意性」は区別すべき概念であり、問題設定に応じて適切なレベルを選ぶ必要がある。

\item \textbf{定義の選択と証明可能性}：最小層という追加構造を導入することで、一般性を保ちつつ完全な普遍性が証明可能となる。これは、適切な定義の選択が定理の証明可能性に決定的な影響を与えることを示している。
\end{itemize}

\subsection{今後の課題}

以下の方向性が今後の研究課題として考えられる：

\begin{enumerate}
\item \textbf{随伴関手の形式化}：忘却関手 $U: \mathbf{StructureTowerWithMin} \to \mathbf{Type}$ と自由構造塔構成 $F$ の間の随伴性を証明する。

\item \textbf{極限・余極限の一般論}：直積以外の圏論的極限（等化子、引き戻しなど）が構造塔の圏で存在するかを調べる。

\item \textbf{Grothendieck トポスへの拡張}：構造塔を層の理論やトポス理論と関連付ける。

\item \textbf{応用}：構造塔の概念を具体的な数学的構造（位相空間、順序構造、代数系）に適用し、Bourbaki の構造主義数学を現代的に再解釈する。
\end{enumerate}

\subsection{謝辞}

本稿は、Claude Code (Anthropic) による \TeX\ 文書生成と、Codex (OpenAI) による Lean 形式化の協働により作成されました。形式的証明の詳細は、添付の Lean ファイル（CAT1.md, CAT2\_advanced.md, CAT2\_revised.md, CAT2\_complete.md）を参照してください。

\begin{thebibliography}{99}
\bibitem{bourbaki}
N. Bourbaki, \textit{Theory of Sets}, Elements of Mathematics, Springer, 1968.

\bibitem{maclane}
S. Mac Lane, \textit{Categories for the Working Mathematician}, 2nd ed., Springer, 1998.

\bibitem{leinster}
T. Leinster, \textit{Basic Category Theory}, Cambridge Studies in Advanced Mathematics, Cambridge University Press, 2014.

\bibitem{mathlib}
The mathlib Community, \textit{The Lean Mathematical Library}, \url{https://github.com/leanprover-community/mathlib4}, 2024.

\bibitem{awodey}
S. Awodey, \textit{Category Theory}, 2nd ed., Oxford Logic Guides, Oxford University Press, 2010.
\end{thebibliography}

\end{document}
